
\subsubsection{Uncertainty Quantification Methods}

Uncertainty quantification in deep learning encompasses multiple approaches including Bayesian neural networks \citep{gal2016dropout}, Monte Carlo Dropout \citep{gal2016dropout}, deep ensembles \citep{lakshminarayanan2017simple}, and evidential deep learning \citep{sensoy2018evidential}. These methods aim to provide predictive uncertainty estimates alongside point predictions.

A limitation of existing work is the absence of validated baselines quantifying the variance from seed initialization alone under deterministic training. Without such baselines, it is unclear what portion of reported uncertainty reflects the method's contribution versus natural initialization stochasticity. This work measures the baseline variance that can serve as a reference for evaluating uncertainty quantification methods.

\subsubsection{Reproducibility and Random Seed Studies}

The deep learning community has documented random seed effects on model performance. Picard et al. (2021) performed an exhaustive search of 10,000 random seeds on CIFAR-10 ResNet-18, finding seed-dependent variance despite deterministic training. Their work demonstrated the existence of seed variance on complex architectures but did not establish protocols for simpler architectures or practical sample sizes.

Zhou et al. (2025) studied random seed effects in language model fine-tuning. Ghasemzadeh et al. (2023) proposed nested k-fold cross-validation to reduce required sample sizes for model selection. While these works acknowledge seed variability, none provide validated protocols for measuring variance in simple neural networks with bootstrap stability analysis.

\subsubsection{Sample Size Theory}

Rajput and Kumar (2023) provided theoretical guidelines for sample size selection in machine learning experiments, recommending $N\geq 30$ for Central Limit Theorem application when effect size $\geq 0.5$ and accuracy $\geq 80$\%. Their criterion establishes power $\geq 0.85$ for detecting meaningful effects but does not empirically test bootstrap stability specifically for neural network variance estimation.

Sluijterman et al. (2023) studied optimal training of mean-variance estimation networks, focusing on loss function design for uncertainty-aware predictions. Their work assumes variance as a learnable prediction target but does not measure classical variance baselines.

This work differs from prior studies by empirically validating the $N\geq 30$ criterion for neural network variance measurement and testing bootstrap stability (confidence interval width $\leq 50$\%). It establishes measurement protocols for simple MLPs before scaling to complex architectures.
